\section{Compact Chern classes}
\label{sec:compact-data}

We first construct many compactly supported line-bundle data on $Y$
whose Chern classes satisfy $e^2=K_Y\cdot e<0$. This equality will make
the weighted index in the next section equal to one. Sums of distinct
canonical basis elements retain the equality, whereas arbitrary
multiples do not. We therefore seek such sums that become trivial
along a fixed seam, so that they can be cut open and extended to $Y$.
Here a datum is a smooth complex line bundle $E$ together with a fixed
nowhere-zero frame outside a compact set, or, equivalently, fixed frames
on both sufficiently distant ends. It has a relative first Chern class
\[
 e=c_1(E,\text{end frames})\in H_c^2(Y;\ZZ).
\]
Its image under the forgetful map
$j:H_c^2(Y;\ZZ)\to H^2(Y;\ZZ)$ is $c_1(E)$.
For compactly supported degree-two classes $a,b$ we use the pairing
\[
 \langle a,b\rangle=\int_Y a\smile j(b),
 \qquad K_Y\cdot a=\int_Y c_1(K_Y)\smile a.
\]
Thus every integral here has a compactly supported factor. In de Rham
notation these integrals are wedge products of closed representatives;
their values are independent of those representatives.

Choose a regular value $c$ of $h$ and put $Z=h^{-1}(c)$. The height is
proper by Lemma~\ref{surface:height}, so $Z$ is a compact two-sided
oriented three-manifold, possibly disconnected. Its image in each $X_n$
is embedded and will again be denoted by $Z$. The quotient map identifies
the cut-open $X_n$ with
\[
 W_n=h^{-1}([c,c+n]),
\]
with its two boundary copies of $Z$ identified by $S=T^n$. Indeed each
$S$-orbit has one representative with $c<h<c+n$, except for the seam
orbits. No connectedness assertion about $Z$ or $W_n$ is needed.

\begin{lemma}[Many compact classes]\label{compact:classes}
There is a constant $C_0\geq0$, independent of $n$, such that for each
$n$ one can choose smooth line-bundle data $E_1,\ldots,E_d$ on $Y$,
supported in the interior of $W_n$, with
\begin{equation}\label{compact:rank-bound}
 d\geq n b_2(X)-C_0.
\end{equation}
They are obtained by cutting line bundles on $X_n$ that are trivialized
near $Z$. Their transferred Chern classes
$w_1,\ldots,w_d\in H^2(X_n;\ZZ)$ are linearly independent over $\mathbb Q$
modulo torsion. Their compact Chern classes $e_i$ satisfy
\begin{equation}\label{compact:numerical-data}
 e_i^2=K_Y\cdot e_i=-m_i<0
\end{equation}
for positive integers $m_i$. Pairings of these compact classes equal
the corresponding pairings of the transferred classes on $X_n$.
\end{lemma}

\begin{proof}
Write $N=n b_2(X)$. By Lemma~\ref{surface:covers}, choose actual integral
lifts $u_1,\ldots,u_N\in H^2(X_n;\ZZ)$ of a canonical negative
orthonormal basis modulo torsion. We first show that every restriction
$u_j|_Z$ is torsion.

Let $i:Z\hookrightarrow X_n$ be the inclusion. For
$a\in H_2(Z;\mathbb Q)$, a small displacement in the normal direction gives
a representative of $i_*a$ disjoint from a representative on $Z$.
This displacement is isotopic to the inclusion, so
$(i_*a)^2=0$. Negative definiteness of $X_n$ forces $i_*a=0$.
In particular the whole image of $H_2(Z;\mathbb Q)$ is zero. By the
cohomology--homology pairing over $\mathbb Q$, the restriction
\[
 H^2(X_n;\mathbb Q)\longrightarrow H^2(Z;\mathbb Q)
\]
is zero. All integral degree-two classes therefore restrict to the fixed
finite group
$G=\operatorname{Tors}H^2(Z;\ZZ)$.

Group the chosen basis lifts by their restriction $\tau\in G$.
Let $k_\tau\geq1$ be the order of $\tau$, with order one for the zero
element, and let $N_\tau$ be the size of its group. Any sum of
$k_\tau$ distinct basis lifts in this group restricts to zero on $Z$.
If $N_\tau>k_\tau$, these sums span the full rational coordinate space
of the group. To check this, take any distinct indices $i,j$ in the
group and choose a set $J$ of $k_\tau-1$ other indices. The difference
\[
 \left(u_i+\sum_{a\in J}u_a\right)
 -\left(u_j+\sum_{a\in J}u_a\right)=u_i-u_j
\]
lies in the span. Such differences span the coordinate-sum-zero
hyperplane, while any one of the indicated sums has nonzero coordinate
sum $k_\tau$. If $N_\tau=k_\tau$, there is one such sum and the loss
of dimension is $k_\tau-1$. If $N_\tau<k_\tau$, taking no sum loses
at most $k_\tau-1$ dimensions. Consequently the collection of all these
sums spans a space of dimension at least
\[
 N-C_0,\qquad C_0=\sum_{\tau\in G}(k_\tau-1).
\]
Choose a rationally independent subcollection of the actual sums, and
call its elements $w_1,\ldots,w_d$. The constant $C_0$ depends only on
the fixed seam $Z$.

A smooth complex line bundle is classified by its integral first Chern
class. Thus $w_i$ is the Chern class of a smooth line bundle $L_i$ on
$X_n$. Since $w_i|_Z=0$ integrally, $L_i$ is trivial on a collar of
$Z$; choose one frame there. Cut along $Z$, retaining the restrictions
of this same frame on the two boundary collars of $W_n$. Lift the
cut-open bundle to $W_n\subset Y$ and extend it over the two remaining
ends using trivial bundles and these frames. This produces $E_i$ with
the stipulated end framings. Its relative class $e_i$ is supported away
from the boundary of $W_n$.

For the pairing assertion, choose a smooth unitary connection on $L_i$
that is trivial in its chosen collar frame. Its normalized curvature is
a closed representative of $w_i$, vanishing near $Z$. The same form,
lifted to the cut-open period and extended by zero, represents $e_i$.
Integrating a wedge product of two such forms on $Y$ is exactly
integrating their original wedge product on $X_n$. The same argument
works with the canonical class, whose curvature form pulls back under
the covering. If $w_i$ is the sum of $m_i$ distinct canonical basis
elements, then
\[
 w_i^2=-m_i,\qquad K_{X_n}\cdot w_i=-m_i.
\]
Torsion parts contribute zero to these integer pairings. This proves
\eqref{compact:numerical-data} and the lemma.
\end{proof}
